% Eksempelfil: vanlige pakker og miljøer i fysikk- og mattedokumenter. Kan ikke kompileres alene.
\documentclass[a4paper,11pt]{article}
\usepackage[margin=2cm, top=2.5cm]{geometry}
\usepackage[per-mode=symbol]{siunitx}
\usepackage[version=4]{mhchem}
\usepackage{tcolorbox,enumitem,circuitikz,pgfplots}
\hypersetup{colorlinks=true, linkcolor=darkred}
\sisetup{output-decimal-marker={,}, separate-uncertainty=true}
\pgfplotsset{compat=1.18}
\graphicspath{{figurer/}}
\addbibresource{kilder.bib}

\begin{document}

\section{Enheter og tall}

Tyngdeakselerasjonen er \qty{9,81}{m/s^2}, og lysfarten er \qty{2,998e8}{\metre\per\second}. Konstanten er \num{6,674e-11}, vinkelen er \ang{30}, og enheten for energi er \unit{\joule}. I matte: $g = \qty{9,81}{m/s^2}$ og $E = \qty{1,6e-19}{\joule}$.

Kjernereaksjonen \ce{^{235}_{92}U + ^1_0n -> ^{141}_{56}Ba + ^{92}_{36}Kr + 3 ^1_0n} frigjør energi, og vann er \ce{H2O}.

\begin{enumerate}[label=\alph*), noitemsep]
    \item Se figur~\ref{fig:krets}, \cref{tab:data} og \autoref{eq:ohm}.
    \item Ifølge \textcite{griffiths2017} og \parencite[s.~12]{einstein1905}.
\end{enumerate}

\begin{tcolorbox}[colback=white, colframe=darkred, title=Ohms lov, boxrule=0.5pt]
    Spenningen over en motstand er $U = RI$.
\end{tcolorbox}

\begin{figure}[htbp]
    \centering
    \includegraphics[width=0.8\linewidth]{figurer/bane.pdf}
    \caption{Banen.}\label{fig:bane}
\end{figure}

\begin{figure}[htbp]
    \centering
    \begin{circuitikz}[scale=0.9]
        \draw (0,0) to[battery1, l=$\varepsilon$] (0,3)
              to[R=$R_1$, i>^=$I$] (4,3)
              to[C, l_=$C$] (4,0) -- (0,0);
    \end{circuitikz}
    \caption{Krets.}\label{fig:krets}
\end{figure}

\begin{tikzpicture}
    \begin{axis}[xlabel={$t$ / \unit{s}}, ylabel={$v$}, ymin=0, width=8cm, grid=major]
        \addplot[darkorange, thick, domain=0:5, samples=50] {9.81*x};
        \addplot+[only marks] coordinates {(0,0) (1,9.8) (2,19.6)};
    \end{axis}
\end{tikzpicture}

\begin{table}[htbp]
    \centering
    \begin{tabular}{@{}l S[table-format=2.2] c@{}}
        \toprule
        Planet & {Avstand (\unit{AU})} & Måner \\
        \midrule
        Jorda & 1,00 & 1 \\
        \bottomrule
    \end{tabular}
    \caption{Data.}\label{tab:data}
\end{table}

\vspace{2mm}\hspace{1em}\noindent Slutt.\\[4pt]
\input{kapittel1}
\printbibliography
\end{document}
